Put
\[
J(s_0,s_1)=\Pr\!\left\{s_1(O_{T,1})\le Q_{k_1}(s_1(O_{1,1}),\ldots,s_1(O_{M,1}))\ \text{and}\ s_0(O_{T,0})\le Q_{k_0}(s_0(O_{1,0}),\ldots,s_0(O_{M,0}))\right\}.
\]
By \(\texttt{def:coverage-frontier}\),
\[
\Gamma_{M,\alpha,\boldsymbol\varepsilon}(w)
=\max_{s_0,s_1}\bigl[(1-\alpha)-J(s_0,s_1)\bigr]_+,
\tag{1}
\]
where the maximum ranges only over deterministic maps \(s_a:\mathcal O\to[0,K]\). Thus (1) is, by construction, the exact worst deficit of the displayed joint rank event in that deterministic orbit-score submodel.

That submodel is strictly narrower than the score behavior allowed by \(\mathcal P_{\mathcal G,K}\). For example, fix the covariates, take both predictors to be zero, draw a nondegenerate Bernoulli variable \(\xi_i\), and set \(Y_{iv}(a)=K\xi_i\) for every member \(v\) and arm \(a\). This law is bounded and invariant under every member relabelling, but, conditional on any selected orbit, its residual score is the nondegenerate random variable \(K\xi_i\). Therefore an arbitrary law in the full model need not induce a deterministic map \(s_a\), and the optimization domain in (1) is not the full-model domain. Hence (1) supplies neither a uniform full-model joint-inclusion certificate nor a causal effect-set coverage certificate.

The event containment in the theorem is nevertheless valid for every realization. Define
\[
E_{\mathrm{joint}}=\{Y_{0T}(1)\in I^{\boldsymbol\varepsilon}_{1,w},\ Y_{0T}(0)\in I^{\boldsymbol\varepsilon}_{0,w}\},
\qquad
E_{\mathrm{eff}}=\{\tau_T\in\mathcal C^{\boldsymbol\varepsilon}_w\}.
\]
On \(E_{\mathrm{joint}}\), the definition of the Minkowski difference and \(\tau_T=Y_{0T}(1)-Y_{0T}(0)\) give
\[
\tau_T\in I^{\boldsymbol\varepsilon}_{1,w}-I^{\boldsymbol\varepsilon}_{0,w}
=\mathcal C^{\boldsymbol\varepsilon}_w.
\tag{2}
\]
Thus \(E_{\mathrm{joint}}\subseteq E_{\mathrm{eff}}\), and monotonicity of conditional probability proves the stated inequality. The reverse containment fails when both intervals are \(\{0\}\) and both potential outcomes equal \(K\): then both armwise inclusions fail, whereas \(\tau_T=0\) belongs to the difference set \(\{0\}\). Equation (2) is algebraic and does not enlarge the optimization domain in (1).

Suppose first that \(k_0,k_1\le M\). For a deterministic score map, the order-statistic identity
\[
s_a(O_{T,a})\le Q_{k_a}(s_a(O_{1,a}),\ldots,s_a(O_{M,a}))
\quad\Longleftrightarrow\quad
\sum_{i=1}^M\mathbf 1\{s_a(O_{i,a})<s_a(O_{T,a})\}<k_a
\tag{3}
\]
shows that the event depends on \(s_a\) only through its total preorder on the finite set \(\mathcal O_a\). Conversely, if the ordered nonempty tie classes are \(C_{a,1}\prec\cdots\prec C_{a,d_a}\), assign score zero when \(d_a=1\), and otherwise assign \((j-1)K/(d_a-1)\) to class \(C_{a,j}\). This realizes the preorder in \([0,K]\) with at most \(m_a\) levels. The maximum in (1) is therefore attained by one of the \(F_{m_0}F_{m_1}\) preorder pairs.

For each such pair, \(\texttt{lem:coherent-two-finite-rank-frontier-evaluator}\) proves
\[
J(C)=\sum_{\pi\in\Pi_C}\omega_C(\pi)D_{M,k_0,k_1}(\pi),
\tag{4}
\]
while retaining the full calibration law \(R_w\), the coherent common-target pushforward \(P_T^C\), and the local weights. That lemma, together with \(\texttt{lem:four-cell-rank-event-dynamic-program}\), \(\texttt{lem:four-cell-multinomial-sum-alternative}\), and \(\texttt{lem:corner-oriented-four-cell-rank-evaluator}\), proves the seven disjoint key strata, every displayed exact evaluator, the branch charge \(\mathsf B\), both enumeration bounds, and memory \(O(m_0m_1+L+U_{\max}+G+\mathsf W)\). It also proves that face and rank-one keys use constant auxiliary workspace and do not trigger factorial preprocessing. The same coherent-evaluator lemma proves the edge, face, and rank-one audit identities, their legal two-block realizations, their displayed level choices, and each strict resource comparison.

If exactly one rank \(k_b=M+1\), \(\texttt{prop:asymmetric-rank-boundary}\) makes the boundary-arm event certain. For \(c=1-b\), conditioning (3) on \(O_{T,c}=O_t\) gives a binomial count with success probability
\[
\theta(O_t)=\sum_{O:s_c(O)<s_c(O_t)}q_{c,O}.
\]
Averaging over the target orbit gives the formula in the statement. The ordered-class scan and exact shorter-tail recurrence in \(\texttt{lem:one-boundary-orbit-frontier-evaluator}\) give the asserted \(F_{m_c}\)-preorder time and constant auxiliary distribution-function memory. If both ranks equal \(M+1\), both rank events are certain and (1) equals zero.

For completeness, the claimed failure of binary attainment follows from a finite exact witness. Use four audited pair blocks with independent fair one-treated--one-control assignment, calibration block masses
\[
q=(13/20,1/20,1/4,1/20),
\]
and common-target block masses
\[
p=(1/20,1/10,1/4,3/5).
\]
Choose a block according to \(q\) and then its unique observed member in the requested arm. This is legal and equivariant on the invariant support. Take arm \(0\) constant and order arm \(1\) as \(1\prec3\prec2\prec4\). At \(M=3\) and \(k_0=k_1=3\), the joint-event failure probability is
\[
\frac14\left(\frac{13}{20}\right)^3
+\frac1{10}\left(\frac{18}{20}\right)^3
+\frac35\left(\frac{19}{20}\right)^3
=\frac{104957}{160000}.
\]
Since \(1-\alpha=1/2\), its four-level objective is \(24957/160000\). For binary high-score sets \(H_0,H_1\), inclusion--exclusion gives failure probability
\[
p(H_0)\{1-q(H_0)\}^3+p(H_1)\{1-q(H_1)\}^3
-p(H_0\cap H_1)\{1-q(H_0)\}^3\{1-q(H_1)\}^3.
\tag{5}
\]
Exact enumeration of the sixteen choices for each high set gives maximum objective \(759109/5000000\). Direct subtraction yields
\[
\frac{24957}{160000}-\frac{759109}{5000000}
=\frac{83189}{20000000}>0,
\]
so binary maps do not attain (1).

If target and calibration orbit supports are disjoint in a finite-rank arm, assign score \(K\) to every target-supported orbit and zero to every calibration-supported orbit in that arm, and use zero in the other arm. The first rank event fails surely and the second holds surely. Hence \(J=0\), so (1) equals \(1-\alpha\), its largest possible value.

Finally, the full-model statements remain separate. Under exact orbit balance, \(\texttt{thm:common-selector-lp}\) proves the Bonferroni coverage guarantee for the ordinary rank construction. Without exact balance, \(\texttt{prop:active-arm-finite-rank-orbit-gap-bound}\) proves the active-arm total-variation lower bound. Neither result invokes (1). Setting \(\varepsilon_0=\varepsilon_1=\alpha/2\) gives \(k_0=k_1=k\) by \(\texttt{lem:armwise-finite-rank-staircase}\), which yields the symmetric deterministic-submodel specialization. Every computational prepass changes only the exact evaluation of (4), so it changes none of the separate active-arm, pair-audit, empty-library, ordinary-rank, or weighted-orbit semantics.